\begin{figure}[t]
    \centering
    \resizebox{\textwidth}{!}{
        \begin{tikzpicture}[
            x=0.52cm, y=0.52cm,
            neuron/.style={circle, draw, minimum size=9pt, inner sep=0pt},
            arr/.style={-{Stealth[length=3pt, width=2.5pt]}},
        ]

        %% ── Layers ─────────────────────────────────────────────────────────────────
        %%  x positions : -9  -5  -2   2   5   9
        %%  y positions :   4   2  -2  -4

        \foreach \lx/\idx in {-9/1, -5/2, -2/3, 2/4, 5/5, 9/6}{
        \foreach \ly/\jdx in {4/1, 2/2, -2/3, -4/4}{
            \node[neuron] (L\idx N\jdx) at (\lx,\ly) {};
        }
        \draw[dotted, thick] (\lx, 0.75) -- (\lx, -0.75);
        }

        %% Layer labels (1 unit above top neuron)
        \node[font=\scriptsize] at (-9,  5.2) {Layer $1$};
        \node[font=\scriptsize] at (-5,  5.2) {Layer $2$};
        \node[font=\scriptsize] at (-2,  5.2) {Layer $l{-}1$};
        \node[font=\scriptsize] at ( 2,  5.2) {Layer $l$};
        \node[font=\scriptsize] at ( 5,  5.2) {Layer $L{-}1$};
        \node[font=\scriptsize] at ( 9,  5.2) {Layer $L$};

        %% ── Fully-connected pairs ───────────────────────────────────────────────────
        \foreach \i in {1,...,4}{
        \foreach \j in {1,...,4}{
            \draw[arr] (L1N\i) -- (L2N\j);
            \draw[arr] (L3N\i) -- (L4N\j);
            \draw[arr] (L5N\i) -- (L6N\j);
        }
        }

        %% ── Horizontal dots between layer groups ────────────────────────────────────
        \draw[dotted, thick] (-4, 0) -- (-3, 0);
        \draw[dotted, thick] ( 3, 0) -- ( 4, 0);

        %% ── Input vector ────────────────────────────────────────────────────────────
        %%  Centre (-12.5, 0), gap = 1.2  →  entries at y = 2.4, 1.2, 0, -1.2, -2.4

        \node (iv0) at (-12.5,  2.4) {$x_0$};
        \node (iv1) at (-12.5,  1.2) {$x_1$};
        \node (iv2) at (-12.5,  0.0) {$\vdots$};
        \node (iv3) at (-12.5, -1.2) {$x_{n_1-1}$};
        \node (iv4) at (-12.5, -2.4) {$x_{n_1}$};

        %% Brackets  (pad = 0.4, bw = 1.0, serif = 0.25)
        \draw (-13.25, 2.8) -- (-13.5, 2.8) -- (-13.5,-2.8) -- (-13.25,-2.8);
        \draw (-11.75, 2.8) -- (-11.5, 2.8) -- (-11.5,-2.8) -- (-11.75,-2.8);

        \node[align=center, font=\scriptsize] at (-12.5, 4.3)
        {Input\\vector $x$};

        %% Arrows: skip the \vdots row (index 2); connect 0,1,3,4 → neurons 1,2,3,4
        \draw[arr] (iv0.east) -- (L1N1);
        \draw[arr] (iv1.east) -- (L1N2);
        \draw[arr] (iv3.east) -- (L1N3);
        \draw[arr] (iv4.east) -- (L1N4);

        %% ── Output vector ───────────────────────────────────────────────────────────
        %%  Centre (12.5, 0)

        \node (ov0) at (12.5,  2.4) {$\hat{y}_0$};
        \node (ov1) at (12.5,  1.2) {$\hat{y}_1$};
        \node (ov2) at (12.5,  0.0) {$\vdots$};
        \node (ov3) at (12.5, -1.2) {$\hat{y}_{n_L-1}$};
        \node (ov4) at (12.5, -2.4) {$\hat{y}_{n_L}$};

        \draw (11.75, 2.8) -- (11.5, 2.8) -- (11.5,-2.8) -- (11.75,-2.8);
        \draw (13.25, 2.8) -- (13.5, 2.8) -- (13.5,-2.8) -- (13.25,-2.8);

        \node[align=center, font=\scriptsize] at (12.5, 4.3)
        {Model's\\prediction $\hat{y}$};

        \draw[arr] (L6N1) -- (ov0.west);
        \draw[arr] (L6N2) -- (ov1.west);
        \draw[arr] (L6N3) -- (ov3.west);
        \draw[arr] (L6N4) -- (ov4.west);

        \end{tikzpicture}
    }
    \caption{Illustration of a Multilayer Perceptron Architecture.}
    \label{fig:general_architecture_multilayer_perceptron}
\end{figure}